% English version of resumen.tex; keep both in sync.
SCDEE is a system for the digital grading of written exams: paper exams are scanned,
assigned to each student and graded on the scanned image with annotations and rubrics. Its
backend, developed in a separate bachelor's thesis, exposes a \acs{rest} \acs{api}
described with OpenAPI, but lacked an interface through which organisation managers,
coordinators, teachers and students could use it.

This work builds that frontend: a single-page application in React that covers the
administration of the organisation and its subjects, exam creation and configuration,
resolution of scanning issues, grading with annotations, publication of grades, review
requests and the student's view of results. The application is integrated with the real
service contract without inventing operations: it calls 92 of its 94 operations, and the
needs the contract does not cover are documented as twelve gaps, each with its evidence and
a proposal for the backend team. An in-browser simulation of the server, validated against
the OpenAPI schemas, makes it possible to develop and test without it.

Development was carried out by two coding agents working in parallel on separate
branches, under the author's supervision. A self-registering module architecture, declared
file ownership and a set of deterministic checks (tests, static analysis, style and
ownership audits) run before every commit allowed both lines of work to be merged without
code conflicts. The merged code runs 179 automated tests that pass. The main limitation is
that it has not been validated against the real server or with users.

Finally, drawing on the recorded process, the work classifies the stages of building a
frontend from a backend as deterministic, reasoning-based or requiring human decision, and
proposes an architecture to automate them.
